% Fragment for the parent article.  Uses its differential-polynomial
% notation A_m, N, D and u_r; replace labels/macros as needed.
\subsection{Two exact coefficient sectors at every degree}
\label{nfsector:section}

Let $\mathcal A_m$ be defined by
\[
 \mathcal A_0=1,\qquad
 \mathcal A_{m+1}=u_0\mathcal A_m-D\mathcal A_m,
 \qquad D u_r=u_{r+1}.
\]
Equivalently its exponential generating series is
\begin{equation}\label{nfsector:Bell}
 \mathcal A(t)=\sum_{m\ge0}\mathcal A_m\frac{t^m}{m!}
 =\exp\left(\sum_{r\ge0}
       \frac{(-1)^r u_r t^{r+1}}{(r+1)!}\right).
\end{equation}
All series in this subsection are formal power series.  Therefore each
coefficient computation is finite.  Write $\mathcal N$ for the classical
functional-monomial projection along the vector space $D\mathbb Q[u_0,u_1,\ldots]$.
It retains powers of $u_0$ and monomials in which the largest positive
derivative index occurs at least twice.  This projection is linear and
need not respect products.

\begin{lemma}[Translation and centered generation]\label{nfsector:Appell}
For an independent constant $c$,
\begin{equation}\label{nfsector:Appellformula}
 \mathcal N(\mathcal A_m)(u_0+c,u_1,\ldots)
 =\sum_{j=0}^m\binom mj c^{m-j}\mathcal N(\mathcal A_j)(u_0,u_1,\ldots).
\end{equation}
Also
\begin{equation}\label{nfsector:centered}
 \mathcal N\mathcal A(t)
 =\mathcal N\exp\left(\sum_{k\ge0}
       \frac{u_{2k}t^{2k+1}}{2^{2k}(2k+1)!}\right).
\end{equation}
\end{lemma}
\begin{proof}
Substitution $u_0\mapsto u_0+c$ commutes with $D$, and preserves the
span of the functional monomials.  Uniqueness of the direct-sum
decomposition therefore implies that it commutes with $\mathcal N$.
Formula~\eqref{nfsector:Bell} acquires the factor $e^{ct}$ under this
substitution, proving \eqref{nfsector:Appellformula}.

For the second assertion, apply the formal translation $e^{tD/2}$
to \eqref{nfsector:Bell}.  Directly summing its exponent gives
\[
 e^{tD/2}\mathcal A(t)
 =\exp\left(\sum_{k\ge0}
       \frac{u_{2k}t^{2k+1}}{2^{2k}(2k+1)!}\right).
\]
Every coefficient of $(e^{tD/2}-1)\mathcal A(t)$ belongs to the image
of $D$, proving the claim.  In particular every surviving monomial
has factor degree congruent to its weight modulo two: the centered
exponent contains only variables of odd weight.
\end{proof}

\begin{theorem}[All polygamma-square coefficients]\label{nfsector:squares}
For every $r\ge1$ and $m\ge2r+2$,
\begin{equation}\label{nfsector:squarecoef}
 [u_0^{m-2r-2}u_r^2]\,\mathcal N(\mathcal A_m)
       =(-1)^r\binom{m}{2r+2}.
\end{equation}
Thus every square term with exactly two factors of positive derivative
order is known at all weights.
\end{theorem}
\begin{proof}
The operator $D$ and the reduction preserve factor degree.  At weight
$2r+2$, the quadratic part of $\mathcal A_{2r+2}$ is
\[
 \frac{(2r+2)!}{2}
 \sum_{a+b=2r}\frac{u_a u_b}{(a+1)!(b+1)!}.
\]
Repeated integration by parts gives
$\mathcal N(u_a u_{2r-a})=(-1)^{a-r}u_r^2$.
The coefficient in question is consequently
\[
 \frac12\sum_{j=1}^{2r+1}
       (-1)^{j-1-r}\binom{2r+2}{j}=(-1)^r,
\]
because the full alternating binomial sum is zero and its two
endpoint terms are both one.  Formula~\eqref{nfsector:Appellformula},
with the original $u_0$ set to zero, now supplies the factor
$\binom{m}{2r+2}$.
\end{proof}

\begin{theorem}[The complete digamma--trigamma sector]
\label{nfsector:trigamma}
After the normal form has been computed, set $u_r=0$ for every $r\ge2$.
Then, for all $m\ge0$,
\begin{equation}\label{nfsector:sector}
 \left.\mathcal N(\mathcal A_m)\right|_{u_2=u_3=\cdots=0}
 =u_0^m-
 \sum_{s=1}^{\lfloor m/4\rfloor}
 \binom{m}{4s}\frac{(4s-1)!!}{4s^2-1}
       u_0^{m-4s}u_1^{2s}.
\end{equation}
Equivalently, put $y=u_1t^2/2$ and define the entire function
$\operatorname{Shi}(y)=\int_0^y\sinh(v)\,dv/v$.  The exponential
generating series of \eqref{nfsector:sector} is
\begin{equation}\label{nfsector:Shi}
 \left.\sum_{m\ge0}\mathcal N(\mathcal A_m)\frac{t^m}{m!}
       \right|_{u_2=u_3=\cdots=0}
 =\frac{e^{u_0t}}2
       \left(\cosh y+\frac{\sinh y}{y}-y\operatorname{Shi}(y)\right),
\end{equation}
where $\sinh y/y$ has its removable value one at zero.
\end{theorem}
\begin{proof}
Let $I$ be the monomial ideal generated by $u_r u_s$ for $r,s\ge2$.
It is stable under $D$.  The stronger property needed here is
\begin{equation}\label{nfsector:ideal}
                   \mathcal N(I)\subseteq I.
\end{equation}
Indeed, an irreducible monomial in $I$ is retained.  If a reducible
monomial in $I$ has a unique largest index $r$, then $r\ge3$.
Replacing $u_r$ by $u_{r-1}$ in its reduction primitive preserves at
least two factors of index at least two.  The primitive and its
derivative therefore remain in $I$.  Induction through the reduction
proves \eqref{nfsector:ideal}.  This argument is essential; stability
under $D$ alone would not imply stability under an arbitrary projection.

Modulo $I$, the Bell series contains at most one factor of derivative
order at least two:
\begin{equation}\label{nfsector:modI}
 \mathcal A(t)\equiv
 e^{u_0t-u_1t^2/2}
 \left(1+\sum_{r\ge2}\frac{(-1)^ru_rt^{r+1}}{(r+1)!}\right)
 \pmod I.
\end{equation}
For integers $l,k\ge0$ and $r\ge2$, integration by parts yields,
modulo the vector subspace $D\mathbb Q[u_0,u_1,\ldots]+I$,
\begin{equation}\label{nfsector:onehigh}
 u_0^l u_1^k u_r\equiv
 \frac{(-1)^{r-1}(l)_{\underline{r-1}}}{k+r-1}
       u_0^{l-r+1}u_1^{k+r}
       \quad(l\ge r-1).
\end{equation}
For $l<r-1$ the class is zero.  Here the underline denotes a falling
factorial.  To prove the formula, for $r\ge3$ differentiate
$u_0^l u_1^k u_{r-1}$: the term containing $u_2u_{r-1}$ belongs
to $I$, leaving a factor $-l$ and reducing $r$ by one.
The final step is
\[
 u_0^l u_1^k u_2\equiv
       -\frac{l}{k+1}u_0^{l-1}u_1^{k+2}.
\]
If there are too few factors of $u_0$, a zero coefficient appears
before this last step.

Let $d_M=[u_1^M]\mathcal N(\mathcal A_{2M})$ for $M\ge2$.
The term without a higher derivative in \eqref{nfsector:modI}
contributes $(2M)!(-1)^M/(2^MM!)$.  For a term with one $u_r$,
the necessary indices are $l=r-1$ and $k=M-r$, where $2\le r\le M$.
Formula~\eqref{nfsector:onehigh} therefore gives
\begin{align*}
 \frac{d_M}{(2M)!}
 &=\frac{(-1)^M}{2^MM!}
   +\frac{(-1)^M}{M-1}
       \sum_{r=2}^{M}
       \frac{(-1)^{r-1}}{(M-r)!2^{M-r}(r+1)!}\\
 &=\frac{(-1)^M\bigl((-1)^{M+1}-1\bigr)}
       {2^{M+1}(M-1)(M+1)!}.
\end{align*}
For the last equality, use the binomial expansion of
$(1/2-1)^{M+1}$ and remove its terms of indices zero, one and two.
Thus $d_M=0$ for odd $M$, while
\[
 d_M=-\frac{(2M-1)!!}{M^2-1}\qquad(M\ge2\text{ even}).
\]
The pure $u_0$ coefficient is one; there is no retained linear $u_1$
term.  The Appell identity now proves \eqref{nfsector:sector}.

Finally, its generating factor is
\[
 1-\sum_{s\ge1}\frac{y^{2s}}{(2s)!((2s)^2-1)}.
\]
Expanding $\cosh y$, $\sinh y/y$ and $y\operatorname{Shi}(y)$
shows that this is the bracket in \eqref{nfsector:Shi} divided by two.
This also proves that the displayed apparent singularity is removable.
\end{proof}

The order of operations in \eqref{nfsector:sector} is part of its
statement: normal reduction precedes specialization.  No assertion is
being made that the normal projection is multiplicative or that one
may impose $u_2=u_3=\cdots=0$ before integration by parts.
On substituting $u_r=\psi^{(r)}(x)$, these formulas give every
digamma--trigamma coefficient and every pure polygamma-square coefficient
in the finite-part moment reductions.  All omitted normal-form terms
in \eqref{nfsector:sector} contain at least two factors of derivative
order at least two.  Their integrals and the exact endpoint functional
must still be retained in the full identity.
